\documentclass[11pt]{amsart}

\usepackage[T1]{fontenc}
\usepackage{lmodern}
\usepackage{microtype}
\usepackage{mathtools,amssymb,amsthm}
\usepackage[hidelinks]{hyperref}

\hypersetup{
  pdftitle={A dyadic obstruction, and K4, against the thirds conjecture for cubic
  graphs under uniform random embedding}
}

\newtheorem{theorem}{Theorem}
\newtheorem{lemma}[theorem]{Lemma}
\newtheorem{corollary}[theorem]{Corollary}

\DeclareMathOperator{\Aut}{Aut}

\title[A dyadic obstruction, and $K_4$, against the thirds conjecture]
{A Dyadic Obstruction, and $K_4$, against the Thirds Conjecture for Cubic
Graphs under Uniform Random Embedding}

\date{}

\subjclass[2020]{05C10, 05C30, 05C80}
\keywords{random embedding, rotation system, signature, singular edge, facial walk,
cubic graph, cycle double cover, counterexample}

\begin{document}

\begin{abstract}
Ghanbari and \v{S}\'{a}mal conjectured that in a random embedding
$(\pi,\lambda)$ of a bridgeless cubic graph $G$ with $m$ edges, the expected
numbers of bad singular, good singular and regular edges are each exactly
$m/3$. The source does not define the underlying probability measure; we settle
the conjecture under two uniform models of ``random embedding'' --- all signed
rotation systems, and the orientable ones --- and we do not settle it under any
non-uniform or limiting measure. (Sampling signatures at one fixed rotation is
not a third model: by the switching lemma below it is another parametrization of
the first, with the same induced distribution of edge classes.) Under the uniform
measure on the $2^{n+m}$ signed rotation systems all three parts fail, in two
independent ways. First, a divisibility
argument: for a connected edge-transitive cubic graph $\Aut(G)$-equivariance
forces each per-edge class probability to be a single dyadic rational
$A_i/2^{n+m}$, so $\mathbb{E}[\#\text{class }i]=m/3$ would require
$3A_i=2^{n+m}$. Second, an exact census of the $2^{10}=1024$ signed rotation
systems of $K_4$, carried out at one fixed rotation and tabulated
signature by signature in the accompanying program, gives
$(\mathbb{E}[\#\mathrm{bad}],\mathbb{E}[\#\mathrm{good}],\mathbb{E}[\#\mathrm{reg}])
=(\tfrac{27}{16},\tfrac{69}{32},\tfrac{69}{32})$
against the conjectured $m/3=2$, so $m/3$ is not a member of that
multiset and no relabelling of the three classes rescues any part. On the
uniform orientable model the numbers change but the three equalities still
fail. What is contradicted is the exact constant $m/3$; the near-uniform split the
source observed is not in dispute, and the conjecture's only consumer in the
source paper keeps its conclusion, which that paper already proves
unconditionally by a different route.
\end{abstract}

\maketitle

\section{The conjecture, verbatim}

Let $G$ be a graph. Following Mohar and Thomassen \cite{MT}, a
\emph{combinatorial embedding} of $G$ is a pair $(\pi,\lambda)$, where $\pi$
assigns to each vertex a cyclic order of the darts at that vertex and
$\lambda:E(G)\to\{+1,-1\}$ is a \emph{signature}. Ghanbari and \v{S}\'{a}mal
\cite[lines 117--120 of \texttt{main.tex}]{GS} classify an edge $e$ of an
embedded graph as follows: $e$ is \emph{singular} if some facial walk $F$
traverses $e$ twice, and \emph{regular} otherwise; a singular edge traversed
twice in the same direction by $F$ is \emph{good singular}, and one traversed
twice in opposite directions is \emph{bad singular}. They then propose the
following \cite[Conjecture~1, \texttt{main.tex} lines 306--313]{GS}, quoted
character for character from the arXiv \LaTeX{} source:

\begingroup\footnotesize
\begin{verbatim}
\begin{conjecture}\label{rnd-emb-conj}
In a random embedding of a bridgeless cubic graph $G$,
    \begin{enumerate}
        \item The expected number of bad singular edges is $\frac{m}{3}$.
        \item The expected number of good singular edges is $\frac{m}{3}$.
        \item The expected number of regular edges is $\frac{m}{3}$.
    \end{enumerate}
\end{conjecture}
\end{verbatim}
\endgroup

\noindent
The sentence immediately preceding it, at \texttt{main.tex} lines 301--304,
is the one that fixes $m$: ``Let $(\pi,\lambda)$ be a
random embedding of a cubic graph $G$. Results of several computer experiments
counting the number of singular edges in a random embedding of a cubic graph
indicate that the expected number of bad singular edges is $\frac{m}{3}$,
where $m$ is the number of edges in $G$.''

That sentence fixes $m$, and it is all the source says about the probability
space: no measure on embeddings is defined there or elsewhere in \cite{GS}, so
what follows depends on a choice of measure that we must make and state. We
work with the uniform measure on signed rotation systems, which is the space
the source's own Mohar--Thomassen notation $(\pi,\lambda)$ names, and we return
in \S\ref{sec:scope} to the uniform orientable model and to what is left
open. Throughout, $n=|V(G)|$, $m=|E(G)|=3n/2$, and
\[
 \Omega=\Omega(G)=\{(\pi,\lambda)\}, \qquad |\Omega|=2^{n}\cdot 2^{m}=2^{n+m},
\]
with the uniform measure: at each vertex of a cubic graph there are exactly
$2$ cyclic orders of the $3$ darts, and $\lambda$ ranges over $\{\pm1\}^{E}$.
Since $n$ is even, $m/3=n/2$ is an integer, so no crude integrality argument is
available.

The quantifier over $G$ is read as universal, i.e.\ one counterexample graph
suffices. That is the source's own use of the statement: its only consumer,
Theorem~4.1 at \texttt{main.tex} lines 317--319, is proved at lines 320--322 by
a first-moment argument applied to the single fixed $G$ it is handed, which
consumes $\mathbb{E}[\cdot]$ for that $G$ and not an average over a family.

Two lines of work are adjacent to Conjecture~1 of \cite{GS} and settle no part
of it. Bender and Richmond \cite{BR} bound the number of singular edges of an
embedding of a $3$-edge-connected graph in terms of the Euler characteristic of
the surface --- a deterministic worst-case bound, $0$ on the sphere and $1$ in
the projective plane --- with no random model and no expectation. The
random-embedding papers of the source's co-author and collaborators --- Campion
Loth and Mohar \cite{CLM}, and Campion Loth, Halasz, Ma\v{s}a\v{r}\'{\i}k,
Mohar and \v{S}\'{a}mal \cite{CLHMMSa,CLHMMSb} --- compute expectations over
uniformly random rotation systems, but of face counts only: none of them
computes the expected number of bad singular, good singular or regular edges,
and none addresses the $m/3$ split. The total-embedding-distribution literature,
for instance Chen, Ou and Zou \cite{COZ}, enumerates signed rotation systems by
genus rather than by edge class, and so does not record the class counts used
below.

\section{A dyadic obstruction}

\begin{theorem}\label{thm:dyadic}
Let $G$ be a finite connected edge-transitive cubic graph, and let $\Omega$
carry the uniform measure. Then for each of the three classes
$i\in\{\mathrm{bad\ singular},\ \mathrm{good\ singular},\ \mathrm{regular}\}$
one has $\mathbb{E}[\#\text{class }i]\neq m/3$; so each of the three parts of
the conjecture fails for $G$ in this model.
\end{theorem}

\begin{proof}
Put $N=|\Omega|=2^{n+m}$.
\begin{enumerate}
\item For each edge $e$ and each class $i$, the number $A_{e,i}$ of
$(\pi,\lambda)\in\Omega$ in which $e$ lies in class $i$ is a non-negative
integer, so $\mathbb{P}[e\in\text{class }i]=A_{e,i}/N$ is \emph{dyadic}.
\item Each $\sigma\in\Aut(G)$ acts on $\Omega$ by
$(\pi,\lambda)\mapsto(\pi^{\sigma},\lambda^{\sigma})$. Being a group action it
is a bijection of the $N$ systems, and it carries facial walks to facial walks
preserving the number of traversals of an edge and their relative directions.
Hence $e$ lies in class $i$ under $(\pi,\lambda)$ if and only if $\sigma(e)$
lies in class $i$ under the image, and $A_{e,i}=A_{\sigma(e),i}$.
\item Edge-transitivity makes all $A_{e,i}$ equal to a common $A_i$, so by
linearity $\mathbb{E}[\#\text{class }i]=\sum_{e}A_{e,i}/N=m\,A_i/N$.
\item If $\mathbb{E}[\#\text{class }i]=m/3$ then $A_i/N=1/3$, i.e.
$3A_i=2^{n+m}$, which is impossible because $3$ divides no power of $2$. This
argument is run for each $i$ separately.\qedhere
\end{enumerate}
\end{proof}

Three remarks fix the strength of Theorem~\ref{thm:dyadic}.

The divisibility must be applied per edge, never to the total: for a cubic
graph $m=3n/2$ is always divisible by $3$, so $m/3=n/2$ is an integer and
``$3$ divides no power of $2$'' says nothing about
$\mathbb{E}[\#\text{class }i]=m/3$ directly. Edge-transitivity is what converts
the total into a per-edge statement, and without it the argument does not run.

``Bridgeless'' is redundant in the hypothesis: a connected edge-transitive
graph with a bridge has every edge a bridge, hence is a tree, and a tree has a
leaf and so is not cubic.

The hypothesis is met by infinitely many graphs --- infinitely many connected
cubic arc-transitive graphs exist, the sextet construction of Biggs and Hoare
\cite{BH} being one source --- but the conclusion is quantitatively weak:
\[
 \bigl|\mathbb{E}[\#\text{class }i]-\tfrac{m}{3}\bigr|
 =\frac{m}{N}\Bigl|A_i-\frac{N}{3}\Bigr|\ \ge\ \frac{m}{3N},
\]
since $A_i$ is an integer and $3\nmid N$; on $K_4$ this reads $1/512$. The floor
improves to $1/32$ once one knows in addition that $2^{n}\mid A_i$, which holds
because the switching action of Lemma~\ref{lem:gauge} is free with all orbits of
size $2^{n}$ and each class is constant on orbits: the nearest multiple of $16$
to $N/3=1024/3$ is $16/3$ away, and $\tfrac{m}{N}\cdot\tfrac{16}{3}=1/32$.
By itself the theorem would license the reply
``the exact equality is impossible, restate it asymptotically''; the
computation of \S\ref{sec:k4} shows the deviation on $K_4$ is ten times that
floor.

\section{The \texorpdfstring{$K_4$}{K4} computation}\label{sec:k4}

\subsection*{Face tracing}
Let $G$ be cubic. Face tracing runs on the $4m$ \emph{arrival states}
$(f,y,t)$, where $f\in E(G)$, $y$ is an endpoint of $f$ and $t\in\{+1,-1\}$,
under $\Phi=T\circ R$:
\begin{itemize}
\item $R$ (rotate at $y$): from $(f,y,t)$, depart $y$ along
$g=\pi_{y}^{\,t}(f)$, keeping $t$;
\item $T$ (traverse $g$ from $y$): arrive at
$\bigl(g,z,t\cdot\lambda(g)\bigr)$, where $z$ is the far end of $g$.
\end{itemize}
The reversal $\operatorname{rev}(f,u,t)=(f,v,-t\lambda(f))$, with $v$ the
other end of $f$, is an involution satisfying
$\operatorname{rev}\circ\Phi\circ\operatorname{rev}=\Phi^{-1}$, so facial
walks correspond to pairs $\{C,\operatorname{rev}(C)\}$ of $\Phi$-cycles and
the number of faces is $f=(\#\Phi\text{-orbits})/2$. Writing $A=(f,u,+1)$ and
$B=(f,u,-1)$ for an edge $f=uv$, the definitions of \cite{GS} read
\begin{equation}\label{eq:class}
\begin{gathered}
 f\ \text{good singular}\iff A\sim_{\Phi}B,\qquad
 f\ \text{bad singular}\iff A\sim_{\Phi}\operatorname{rev}(B),\\
 \text{and $f$ regular otherwise},
\end{gathered}
\end{equation}
and these are mutually exclusive and exhaustive (both would force
$A\sim_{\Phi}\operatorname{rev}(A)$, which is impossible). The conjecture's
``$1/3$ each'' shape is the guess that the three possible same-cycle relations
among an edge's four tracing states are equally likely.

\subsection*{The switching lemma}
Switching at a vertex $v$ means replacing $\pi_v$ by $\pi_v^{-1}$ \emph{and}
flipping $\lambda$ on the three edges at $v$; this leaves the facial walks,
hence all three edge classes, unchanged \cite[\texttt{main.tex} lines
158--161]{GS}. Switching at $S\subseteq V$ therefore flips $\pi$ on $S$ and
flips $\lambda$ on the cut $\delta(S)$, and for fixed $S$ the map
$\lambda\mapsto\lambda\cdot\delta(S)$ is a bijection of $\{\pm1\}^{E}$.

\begin{lemma}\label{lem:gauge}
$\mathbb{Z}_2^{V}$ acts on $\Omega$ by switching, the action is free with all
orbits of size $2^{n}$, each orbit meets $\{\pi=\pi_0\}$ in exactly one point,
and all three edge classes are constant on orbits. Consequently, for every
fixed rotation $\pi_0$ and every class $i$,
\[
 \mathbb{E}_{(\pi,\lambda)\sim\Omega}\bigl[\#\text{class }i\bigr]
 \;=\;
 \mathbb{E}_{\lambda\sim\{\pm1\}^{E}}\bigl[\#\text{class }i \,\bigm|\, \pi=\pi_0\bigr]
 \qquad\text{exactly.}
\]
\end{lemma}

\begin{proof}
Switching at $S$ changes $\pi$ at exactly the vertices of $S$, so the
stabiliser of any point is trivial and orbits have size $2^{n}$; and the orbit
of $(\pi_0,\lambda)$ meets $\{\pi=\pi_0\}$ only at $S=\emptyset$. Class
invariance is the vertex statement iterated over $S$. Averaging a
class-constant function over $\Omega$ is therefore averaging it over the $2^m$
orbit representatives $(\pi_0,\lambda)$.
\end{proof}

\subsection*{The count}
Take $K_4$ on $\{0,1,2,3\}$ with the six edges indexed $0,\dots,5$ in the
order
\[
 (0,1),\ (0,2),\ (0,3),\ (1,2),\ (1,3),\ (2,3),
\]
so that dart $2i$ is the end of edge $i$ at its smaller vertex and dart $2i+1$
the other end. Let $\pi_0$ take, at each vertex, the three incident darts in
increasing dart order; explicitly
\[
 \pi_0: \quad 0\mapsto(1,2,3),\quad 1\mapsto(0,2,3),\quad
 2\mapsto(0,1,3),\quad 3\mapsto(0,1,2)
\]
as cyclic orders of neighbours. Tracing \eqref{eq:class} at $(\pi_0,\lambda)$
for each of the $2^{6}=64$ signatures $\lambda$ classifies all six edges; the
per-signature table --- signature, number of faces, Euler genus $2-(n-m+f)$,
the six classes and the triple $(\text{bad},\text{good},\text{reg})$ --- is
recorded in the accompanying program \texttt{verify.py} rather than printed
here, and any single row of it can be redone by hand. Over the $64$ signatures
the three class columns sum to $108$, $138$ and $138$.


\begin{theorem}\label{thm:k4}
For $G=K_4$, which is edge-transitive, every edge is bad singular in exactly
$288$ of the $1024$ signed rotation systems, good singular in exactly $368$,
and regular in exactly $368$. Hence
\[
 \mathbb{P}[e\ \mathrm{bad}]=\tfrac{288}{1024}=\tfrac{9}{32},\qquad
 \mathbb{P}[e\ \mathrm{good}]=\mathbb{P}[e\ \mathrm{regular}]
 =\tfrac{368}{1024}=\tfrac{23}{64},
\]
and by linearity over the six edges
\[
 \mathbb{E}[\#\mathrm{bad}]=\tfrac{27}{16}=1.6875,\qquad
 \mathbb{E}[\#\mathrm{good}]=\mathbb{E}[\#\mathrm{regular}]
 =\tfrac{69}{32}=2.15625,
\]
against the conjectured value $m/3=2$; in particular
$2\notin\{\tfrac{27}{16},\tfrac{69}{32},\tfrac{69}{32}\}$.
\end{theorem}

\begin{proof}
By Lemma~\ref{lem:gauge} it suffices to average over the $2^6=64$ signatures
at the single rotation $\pi_0$; each is classified by \eqref{eq:class} from a
trace of the $24$ arrival states of $K_4$, tabulated row by row in
\texttt{verify.py}. The class column
sums are $108$, $138$, $138$, so by Lemma~\ref{lem:gauge} the number of
(system, edge) pairs over all of $\Omega$ with the edge bad singular is
$\tfrac{1024}{64}\cdot 108=1728$, and $\tfrac{1024}{64}\cdot 138=2208$ for each
of the other two classes. Now $\Aut(K_4)=S_4$ is transitive on the six edges,
so by step~(2) of the proof of Theorem~\ref{thm:dyadic} the count $A_{e,i}$ is
independent of $e$; dividing by $6$ gives $288$, $368$, $368$ per edge. The
probabilities and, by linearity, the expectations follow.
\end{proof}

\begin{corollary}\label{cor:k4}
In the uniform model $\Omega$, the three expectations of $K_4$ are all
different from $m/3$, and no relabelling of the three classes changes that.
Since $K_4$ is a simple, $3$-connected, bridgeless cubic graph, each of the
three parts of the conjecture fails on it under this measure.
\end{corollary}

\begin{proof}
The conjecture asserts that all three class expectations equal $m/3$, so a
relabelling permutes parts (1), (2), (3) and leaves the asserted set of
equalities unchanged; by Theorem~\ref{thm:k4} the value $m/3=2$ is not an
element of the multiset of the three expectations.
\end{proof}

\subsection*{The orientable signatures}
The system $(\pi_0,\lambda)$ is orientable exactly when $\lambda$ is a
coboundary $\delta(S)$. The $16$ vertex subsets of $K_4$ collapse $2$-to-$1$
onto $8$ distinct signatures, namely
\begingroup\small
\begin{verbatim}
   ++++++    ---+++    -++--+    +-+-+-
   ++-+--    +----+    -+--+-    --++--
\end{verbatim}
\endgroup
\noindent
(in the same edge order as above). Their bad counts are
$2,3,3,3,3,2,0,2$, with \emph{good count $0$
throughout}, exactly as \cite[\texttt{main.tex} line 119]{GS} demands
(``in orientable embeddings all singular edges are bad singular''); their
Euler genera are all even, as orientability also demands. Neither fact was
used in the tracing, so both test the good/bad convention against the
source's own assertion.

\section{A sign-flip involution}

Line 304 of \cite{GS} says the expectation ``seems to be the same for good
singular edges and regular edges''. It is the same, exactly, for every edge of
every loopless graph.

\begin{theorem}\label{thm:involution}
Let $G$ be a graph without loops and $e\in E(G)$. Then flipping $\lambda(e)$ is
a measure-preserving involution of $\Omega$ that exchanges the event
$\{e\ \mathrm{regular}\}$ with $\{e\ \mathrm{good\ singular}\}$ and fixes
$\{e\ \mathrm{bad\ singular}\}$; and it changes the number of faces by $-1$,
$+1$ and $0$ respectively. Hence
\begin{equation}\label{eq:involution}
 \mathbb{P}[e\ \mathrm{good}]=\mathbb{P}[e\ \mathrm{regular}]
 \quad\text{exactly},\qquad
 \mathbb{E}[\#\mathrm{bad}]=m-2\,\mathbb{E}[\#\mathrm{good}],
\end{equation}
so parts (2) and (3) of the conjecture are each equivalent to part (1), which
is the aggregate assertion $\mathbb{E}[\#\mathrm{bad}]=m/3$: equivalently, that
a uniformly random edge of $G$ is bad singular with probability $1/3$. No
equality of the individual per-edge probabilities is asserted.
\end{theorem}

\begin{proof}
Write $e=uv$, $\lambda'=\lambda\cdot\mathbf{1}_{e}$, and let $\Phi,\Phi'$ be
the tracing permutations of $(\pi,\lambda)$ and $(\pi,\lambda')$. Since
$\lambda\mapsto\lambda'$ is a bijection of $\{\pm1\}^{E}$ it is measure
preserving, so only the classification has to be tracked.

Let $S$ be the four states on $e$, and put $A=(e,u,+1)$, $B=(e,u,-1)$,
$P=\operatorname{rev}(A)$, $Q=\operatorname{rev}(B)$, so that
$\{P,Q\}=\{(e,v,+1),(e,v,-1)\}$ and $\operatorname{rev}|_{S}=(A\,P)(B\,Q)$.
Let $\tau=(A\,B)(P\,Q)$ be the involution of $S$ negating the sign bit.
Because $e$ is not a loop, $\Phi$ leaves $S$ at the first step, and a step of
$\Phi$ \emph{arrives} in $S$ exactly when it traverses $e$; that step is the
only place $\lambda(e)$ is consulted. Hence $\Phi'$ agrees with $\Phi$ off
those steps and satisfies $\Phi'=\tau\circ\Phi$ there, so: every $\Phi$-orbit
disjoint from $S$ is also a $\Phi'$-orbit, and the first-return maps
$\sigma,\sigma':S\to S$ of $\Phi,\Phi'$ satisfy
\begin{equation}\label{eq:return}
 \sigma'=\tau\circ\sigma .
\end{equation}
Moreover $\#\Phi\text{-orbits}=\#\sigma\text{-cycles}+K$, where $K$ counts the
$\Phi$-orbits disjoint from $S$, and the same with primes and the same $K$.
Two states of $S$ are $\Phi$-equivalent iff they lie in one $\sigma$-cycle, so
by \eqref{eq:class}, $e$ is good singular iff $A,B$ lie in one $\sigma$-cycle
and bad singular iff $A,Q$ do; and under $\lambda'$ one has
$\operatorname{rev}'(A)=Q$, $\operatorname{rev}'(B)=P$, so $e$ is good
singular iff $A,B$ lie in one $\sigma'$-cycle and bad singular iff $A,P$ do.

Now $\operatorname{rev}\circ\Phi\circ\operatorname{rev}=\Phi^{-1}$ and
$\operatorname{rev}(S)=S$ give
$\operatorname{rev}\sigma\operatorname{rev}=\sigma^{-1}$ on $S$; and no
$\sigma$-cycle contains both $A$ and $P=\operatorname{rev}(A)$, nor both $B$
and $Q$, since a $\Phi$-cycle is never its own mirror. Enumerating the
$\sigma$-cycle containing $A$ under those two constraints leaves exactly three
possibilities, and \eqref{eq:return} determines each image:
\[
\begin{array}{lll}
 \sigma=\mathrm{id}       & \sigma'=(A\,B)(P\,Q) & \text{regular}\to\text{good},\\
 \sigma=(A\,B)(P\,Q)     & \sigma'=\mathrm{id}   & \text{good}\to\text{regular},\\
 \sigma=(A\,Q)(B\,P)     & \sigma'=(A\,P)(B\,Q) & \text{bad}\to\text{bad}.
\end{array}
\]
(For the first case: if $A$ is a fixed point of $\sigma$ then
$\sigma^{-1}(P)=\operatorname{rev}\sigma\operatorname{rev}(P)
=\operatorname{rev}\sigma(A)=P$ and likewise for $Q$, and $B$ cannot be
$\sigma$-joined to $P$ without violating $\sigma^{-1}(P)=P$; the other two
cases are forced the same way from $\sigma(A)=B$ and $\sigma(A)=Q$.) This is
the asserted exchange, and it also proves the trichotomy. Counting
$\sigma$-cycles gives $4,2,2$ respectively, so
$f=(\#\sigma\text{-cycles}+K)/2$ changes by $-1$, $+1$, $0$. The displayed
identities follow from
$\mathbb{P}[\mathrm{bad}]+\mathbb{P}[\mathrm{good}]+\mathbb{P}[\mathrm{reg}]=1$
and linearity over the $m$ edges.
\end{proof}

Neither Theorem~\ref{thm:dyadic} nor \S\ref{sec:k4} uses
\eqref{eq:involution}; it is not needed to settle the conjecture, and is
recorded only because it settles, as a theorem, a coincidence the source
noticed empirically.

\section{What is, and is not, settled}\label{sec:scope}

\subsection*{The source paper keeps its theorem}
The conjecture's only consumer in \cite{GS} is Theorem~4.1
(\texttt{main.tex} lines 317--319): ``Suppose that Conjecture~1 is true and
let $G$ be a bridgeless cubic graph. Then, there exists an embedding of $G$
with at most $\frac{m}{3}$ singular edges.'' The failure of the conjecture in
the models treated here removes that hypothesis, so the proof at lines
320--322 no longer supports the conclusion. The conclusion itself is unaffected, and \cite{GS} already contains
an unconditional proof of it in the referee remark at lines 324--325: take any
perfect matching and extend the complementary $2$-factor to a facial double
cover; then only matching edges can be singular. A bridgeless cubic graph has
a perfect matching \cite{Pet}, and a perfect matching of a cubic graph on $n$
vertices has $n/2=m/3$ edges, so the bound follows with no probabilistic
input. What is lost is a conjecture and one redundant conditional proof, not a
theorem.

\subsection*{The measure is not defined by the source}
``Let $(\pi,\lambda)$ be a random embedding'' is the whole of it, so everything
above depends on our choice of $\Omega$ and the conjecture is settled only
under the models we name. Two uniform models were checked, and the three
equalities fail under each: (i) uniform on the $2^{n+m}$ signed
rotation systems, used above, and the space the source's own Mohar--Thomassen
notation names --- sampling $\lambda$ uniformly at a fixed $\pi$ is not a
further model but the same one in other coordinates, with the same induced
distribution of edge classes and hence identical numbers, by
Lemma~\ref{lem:gauge}; and (ii) uniform on the \emph{orientable}
embeddings, i.e.\ on the $2^{n}$ pure rotation systems, where on $K_4$ the
eight coboundary signatures above give
\[
 \mathbb{E}[\#\mathrm{bad}]=\tfrac{18}{8}=\tfrac94,\quad
 \mathbb{E}[\#\mathrm{good}]=0,\quad
 \mathbb{E}[\#\mathrm{regular}]=\tfrac{15}{4},
\]
against $m/3=2$: parts (1) and (3) fail by $+\tfrac14$ and $+\tfrac74$, and
part (2) fails degenerately by line 119. The sign of the deviation flips
between (i) and (ii), so any statement of the form
``$\mathbb{E}[\#\mathrm{bad}]\le m/3$'' must name its measure. A non-uniform or
limiting measure is not covered: the source reports ``several computer
experiments'', it does not say what they sampled, and if they sampled some
other space then what is settled here is the conjecture under each of the two
uniform models above, not the statement about whatever was sampled there.

\subsection*{Only the exact constant is contradicted}
The computed per-edge probabilities are all near $1/3$ --- on $K_4$,
$\tfrac{9}{32}=0.28125$ for bad and $\tfrac{23}{64}=0.359375$ for good and for
regular --- and separating $\mathbb{E}[\#\mathrm{bad}]$ from $m/3$ by sampling
needs many samples. So what is contradicted is the exact constant $m/3$, not
the near-uniform split the source observed. Whether
$\mathbb{E}[\#\mathrm{bad}]/m$ tends to $1/3$ along some family, and whether
$\mathbb{E}[\#\mathrm{bad}]<m/3$ always holds in the uniform signed-rotation
model, are not addressed here; the latter is in any case false without the
model qualifier, since the orientable model gives $9/4>2$ on $K_4$.

\subsection*{Further limits}
Theorem~\ref{thm:dyadic} covers only edge-transitive cubic graphs; for the
rest only exhaustive counting is offered, on the seven graphs listed in
\S\ref{sec:verification}. Nothing here proves the conjecture fails for
\emph{every} bridgeless cubic graph, though being universally quantified it
needs only one counterexample. The infinitude of the connected cubic
arc-transitive family is quoted from \cite{BH} and standard census material
rather than verified here. The unproved Properties 6--9 of \cite{GS} are used
nowhere above. The identification of the three computed expectations with the
three parts of the conjecture is a reading of the source's prose, made here and
not verified by the program; nor does the program check that $K_4$ as it
hard-codes it is simple and $3$-connected, which is left to the reader.
Finally, the program checks two conditions transcribed from the literature and
used nowhere above: a bound of Bender and Richmond \cite{BR} on the number of
singular edges of an embedding of a $3$-edge-connected graph, in the form the
authors of \cite{GS} themselves state it (their companion e-print
arXiv:2511.07285, \texttt{CDC.tex} line 133), and the orientable genus
distribution of $K_4$ given by Gross and Furst \cite{GF}. Neither the 1990 nor
the 1987 paper was consulted here beyond its published abstract, so those two
checks test the tracer against restatements only, and no claim about either is
made.

\section{Verification}\label{sec:verification}

The accompanying program \texttt{verify.py} (Python 3.9+, standard library
only, exact integer and \texttt{Fraction} arithmetic, no floating-point
decision) implements the face tracing of \eqref{eq:class} from the definitions
and re-derives every number above. It carries the $64$-row table of $K_4$
signature data \emph{as data} and compares it, row by row --- signature,
faces, Euler genus, six-letter class string and triple --- with its own trace
of the $24$ arrival states of $K_4$; it then exhausts the full $2^{n+m}$ space
of $\theta$, $K_4$, $K_{3,3}$ and the $3$-prism, and all $2^{m}$ signatures at
each of two fixed rotations of $Q_3$, Wagner $V_8$ and Petersen, the latter
licensed by Lemma~\ref{lem:gauge}; on each of those seven cells the three
computed expectations all differ from $m/3$. That lemma is tested rather than
assumed:
the program switches at every vertex of every embedding of $\theta$ and $K_4$,
and at every vertex of a bounded sample on the other five cells, and it also
checks directly that the per-edge counters are the same at every rotation it
traces. Theorem~\ref{thm:involution} is likewise tested in its face-count form,
on every (embedding, edge) pair of $\theta$ and $K_4$. The program also
checks the source's line 119 on every orientable embedding, checks
that some embedding of $K_4$ does carry a good singular edge, so that the
negative results are not vacuous, and reproduces row by row the two
transcribed literature conditions recorded in \S\ref{sec:scope}, which are
used nowhere in this paper. Here $\theta$ is the $3$-dipole, a
multigraph, and is a cell of this conjecture only if multigraphs count as
bridgeless cubic, a reading the source neither grants nor denies; nothing above
depends on it. The recorded run of the program is \texttt{verify.output.txt}.

\begin{thebibliography}{9}

\bibitem{BR}
E.~A. Bender and L.~B. Richmond,
\emph{$3$-edge-connected embeddings have few singular edges},
J. Graph Theory \textbf{14} (1990), no.~4, 475--477.
\href{https://doi.org/10.1002/jgt.3190140411}{doi:10.1002/jgt.3190140411}.

\bibitem{CLM}
J.~Campion Loth and B.~Mohar,
\emph{Expected number of faces in a random embedding of any graph is at most
linear},
Combin. Probab. Comput. \textbf{32} (2023).
\href{https://doi.org/10.1017/s096354832300010x}{doi:10.1017/s096354832300010x};
arXiv:2202.07746.

\bibitem{CLHMMSa}
J.~Campion Loth, K.~Halasz, T.~Ma\v{s}a\v{r}\'{\i}k, B.~Mohar and
R.~\v{S}\'{a}mal,
\emph{Random $2$-cell embeddings of multistars},
Proc. Amer. Math. Soc. \textbf{150} (2022), no.~9.
arXiv:2103.05036.

\bibitem{CLHMMSb}
J.~Campion Loth, K.~Halasz, T.~Ma\v{s}a\v{r}\'{\i}k, B.~Mohar and
R.~\v{S}\'{a}mal,
\emph{Random embeddings of graphs: the expected number of faces in most graphs
is logarithmic},
Proc. 35th ACM--SIAM Symposium on Discrete Algorithms (SODA 2024).
arXiv:2211.01032.

\bibitem{COZ}
Y.~Chen, L.~Ou and Q.~Zou,
\emph{Total embedding distributions of Ringel ladders},
arXiv:1011.3869, 2010.

\bibitem{BH}
N.~L. Biggs and M.~J. Hoare,
\emph{The sextet construction for cubic graphs},
Combinatorica \textbf{3} (1983), 153--165.
\href{https://doi.org/10.1007/BF02579289}{doi:10.1007/BF02579289}.

\bibitem{GS}
B.~Ghanbari and R.~\v{S}\'{a}mal,
\emph{Facial diagrams and cycle double cover},
arXiv:2605.01410v1, 2026 (EUROCOMB'25 extended abstract).
\href{https://arxiv.org/abs/2605.01410}{arXiv:2605.01410}.
Line numbers refer to \texttt{main.tex} in the arXiv e-print source, the only
version published.

\bibitem{GF}
J.~L. Gross and M.~L. Furst,
\emph{Hierarchy for imbedding-distribution invariants of a graph},
J. Graph Theory \textbf{11} (1987), no.~2, 205--220.
\href{https://doi.org/10.1002/jgt.3190110211}{doi:10.1002/jgt.3190110211}.

\bibitem{MT}
B.~Mohar and C.~Thomassen,
\emph{Graphs on Surfaces},
Johns Hopkins University Press, Baltimore, 2001.

\bibitem{Pet}
J.~Petersen,
\emph{Die Theorie der regul\"{a}ren Graphs},
Acta Math. \textbf{15} (1891), 193--220.
\href{https://doi.org/10.1007/BF02392606}{doi:10.1007/BF02392606}.

\end{thebibliography}

\end{document}
